\chapter{Familia}

\pantalla{Familia} es el catálogo del local, para consultar sin cambiar nada. Sirve para que todo el personal sepa qué hay, cuánto cuesta y qué lleva cada producto.

\captura{familia}{Catálogo de productos}

\begin{itemize}
  \item Arriba se elige entre \boton{Productos} y \boton{Promociones}.
  \item La búsqueda encuentra productos aunque se escriban sin tildes.
  \item Los botones de categoría filtran el catálogo, y dentro de cada categoría se puede elegir una subcategoría.
  \item Las tarjetas avisan el stock: \textbf{Sin stock} si no alcanza para ninguna porción y \textbf{Agotado} si administración lo marcó como no disponible.
\end{itemize}

\captura[0.8]{familia-detalle}{Detalle de un producto con su receta}

Al tocar un producto se ve su precio, su descripción y su \textbf{receta}: cuánto lleva de cada insumo por porción y para cuántas porciones alcanza el stock de cada uno. Si no alcanza para ninguna, el aviso dice qué insumo falta; el producto se puede vender igual, pero conviene avisar a cocina.

En \boton{Promociones} cada combo muestra qué productos incluye, cuánto se ahorra y cuándo aplica (fechas, días y horario).
